\section{Methods: IBD Detection via Hidden Markov Model}
\label{sec:methods_ibd}

\subsection{Model Formulation}

We formulate IBD detection as a two-state Hidden Markov Model. Given windowed sequence identity values from pangenome alignments, the model distinguishes regions of shared recent ancestry (IBD, state~1) from background similarity (non-IBD, state~0).

\paragraph{Emissions.}
Each state emits identity values $o_t \in [0,1]$ according to a Gaussian distribution $P(o_t \mid s_t = s) = \mathcal{N}(o_t; \mu_s, \sigma_s^2)$. The non-IBD mean derives from coalescent theory: $\mu_0 = 1 - \theta$, where $\theta = 4\mu N_e$ is the population-scaled nucleotide diversity. Population-specific $\pi$ values range from 0.00080 (EAS) to 0.00125 (AFR), yielding expected non-IBD identity of 0.99875--0.99920 (Table~\ref{tab:pop_params}). The IBD mean is $\mu_1 = 0.9997$ with $\sigma_1 = 0.0005$, reflecting only sequencing and assembly errors.

\begin{table}[ht]
\centering
\caption{Population-specific nucleotide diversity and expected sequence identity parameters.}
\label{tab:pop_params}
\begin{tabular}{lccc}
\hline
\textbf{Population} & $\boldsymbol{\pi}$ \textbf{(SNPs/bp)} & $\boldsymbol{E[\text{identity} \mid \text{non-IBD}]}$ & \textbf{SNR} \\
\hline
AFR & 0.00125 & 0.99875 & 1.06 \\
EUR & 0.00085 & 0.99915 & 0.74 \\
EAS & 0.00080 & 0.99920 & 0.69 \\
CSA & 0.00095 & 0.99905 & 0.81 \\
AMR & 0.00100 & 0.99900 & 0.85 \\
\hline
\end{tabular}
\end{table}

The per-window signal-to-noise ratio SNR$_1 = \Delta\mu / \sqrt{\sigma_0^2 + \sigma_1^2}$ ranges from 0.69 (EAS) to 1.06 (AFR). Reliable detection requires accumulating evidence over $L_{\min} = (3/\text{SNR}_1)^2$ consecutive windows, corresponding to minimum detectable segments of $\sim$90\,kb (AFR) to $\sim$190\,kb (EAS). Notably, African populations have the best detection power because the identity gap ($\Delta\mu \approx \theta$) is proportionally larger.

\paragraph{Transitions.}
The transition matrix is parameterized by expected IBD segment length $L$ (in windows) and entry probability $p_{\text{enter}}$: $P(\text{stay IBD}) = 1 - 1/L$, with population-adaptive scaling reflecting coalescent depth. Distance-dependent transitions ($P(\text{switch} \mid d) = 1 - e^{-\lambda d}$) and recombination-rate-aware transitions via genetic maps improve biological realism (Supplementary Section~L.2).

\subsection{Parameter Estimation}

Emission parameters are estimated hierarchically: (1)~k-means clustering into two groups, (2)~EM fallback with MAP regularization if clusters are degenerate ($\Delta\mu \leq 0.0005$), and (3)~BIC model selection to avoid overfitting a two-component model when the data does not support it. All parameters are then refined via Baum-Welch~\cite{baum1970maximization}, which re-estimates emissions and transitions jointly over 20 iterations (convergence details in Supplementary Section~L.3). Biological bounds are enforced after each M-step to maintain identifiability ($\mu_0 < \mu_1$) and numerical stability.

\subsection{Inference and Segment Extraction}

The most likely state sequence is computed via the Viterbi algorithm in log-space; posterior state probabilities $\gamma_t(s)$ are computed via forward-backward. IBD segments are extracted by run-length encoding of the Viterbi path, with boundaries refined using posteriors (extension at $\gamma > 0.5$, trimming at $\gamma < 0.2$). Each segment receives a LOD score $\text{LOD} = \sum_{t} \log_{10} [P(o_t \mid \text{IBD}) / P(o_t \mid \text{non-IBD})]$ following the hap-ibd convention~\cite{browning2020hap_ibd}, and a composite quality score $Q \in [0, 100]$ combining posterior strength, consistency, LOD evidence, and segment length.

An identity floor parameter (default 0.9) filters out alignment-gap windows (identity near zero) before HMM inference, treating them as missing data rather than evidence of non-IBD. This reduces variance from $\sigma \approx 0.29$ (bimodal raw distribution) to $\sigma \approx 0.001$ (aligned regions only). Details of the quality score, identity floor analysis, and batch processing are in Supplementary Section~L.2.
